\documentclass{article}

% NeurIPS 2026 style
\usepackage{neurips_2026}

% Required packages
\usepackage[utf8]{inputenc}
\usepackage[T1]{fontenc}
\usepackage{hyperref}
\usepackage{url}
\usepackage{booktabs}
\usepackage{amsfonts}
\usepackage{amsmath}
\usepackage{amssymb}
\usepackage{graphicx}
\usepackage{multirow}
\usepackage{subcaption}
\usepackage{xcolor}
\usepackage{algorithm}
\usepackage{algorithmic}
\usepackage{wrapfig}
\usepackage{enumitem}
\usepackage{nicefrac}
\usepackage{microtype}
\usepackage{caption}
\usepackage{adjustbox}

% Colors for tables
\definecolor{bestcolor}{HTML}{E8F5E9}
\definecolor{ourscolor}{HTML}{FFF3E0}

% Custom commands
\newcommand{\methodname}{DynaCLIP}
\newcommand{\ie}{\textit{i.e.}}
\newcommand{\eg}{\textit{e.g.}}
\newcommand{\etal}{\textit{et al.}}
\newcommand{\wrt}{\textit{w.r.t.}}
\newcommand{\reals}{\mathbb{R}}
\newcommand{\best}[1]{\textbf{#1}}
\newcommand{\second}[1]{\underline{#1}}
\newcommand{\tbd}[1]{#1\textsuperscript{\dag}}

\title{\methodname{}: Physics-Aware Visual Representations via \\ Category-Grounded Contrastive Pre-Training}

\author{
  Anonymous Authors \\
  Anonymous Institution \\
  \texttt{anonymous@institution.edu}
}

\begin{document}

\maketitle

% ======================================================================
\begin{abstract}
% ======================================================================
Visual representations for robotic manipulation must capture not only \emph{what} an object is and \emph{where} it appears, but \emph{how} it will respond to applied forces---its mass, surface friction, and elasticity. Current visual encoders~\cite{oquab2024dinov2,radford2021clip,zhai2023siglip}, including those specifically designed for robotics~\cite{nair2022r3m,radosavovic2023vc1,karamcheti2023voltron}, are trained on semantic, geometric, or temporal objectives that discard physics-relevant cues. We introduce \methodname{}, a contrastive pre-training framework that imbues visual backbones with sensitivity to latent physical material properties. Our key insight is that \emph{object categories carry strong implicit priors over material types}: a hammer is metal, a teddy bear is fabric. By mapping 345 object categories to 10 material types with physically plausible parameter distributions, we generate category-grounded physics configurations without real-world measurement. We simulate analytical dynamics for each configuration, compute pairwise dynamics similarity, and train a DINOv2-ViT-B/14 backbone with a \emph{Soft InfoNCE} loss using continuous similarity as soft contrastive targets. On the \textbf{DynaCLIP-Bench} suite of 10 experiments, \methodname{} (88.2M params) outperforms frozen DINOv2-L/14 (304M) by +25\% on mass $R^2$ and is the \emph{only} model achieving positive mass $R^2$ on the invisible physics test. On \textbf{LIBERO-10}---an established multi-task manipulation benchmark---\methodname{} achieves \best{59.0\%} average success rate, outperforming DINOv2-B/14 (51.5\%), CLIP-L/14 (46.9\%), VC-1 (31.4\%), and 6 other baselines. \methodname{} further produces the most physics-robust downstream policies (generalization ratio $1.02\times$ vs.\ $0.76\times$ for CLIP) on ManiSkill3. Code, data, and checkpoints will be released upon publication.
\end{abstract}

% ======================================================================
\section{Introduction}
\label{sec:intro}
% ======================================================================

Deploying robots in unstructured environments demands visual representations that go beyond recognizing \emph{what} an object is or \emph{where} it appears. A manipulation policy tasked with grasping a ceramic mug must produce markedly different forces, trajectories, and grasp widths than one handling a foam cup---even though both objects share the semantic label ``cup.'' The distinction lies in their physical material properties: mass, surface friction, and elastic restitution, none of which are captured by existing large-scale visual encoders.

The dominant paradigm for visual backbones in robotics has followed two parallel tracks. \emph{General-purpose} encoders---CLIP~\cite{radford2021clip}, DINOv2~\cite{oquab2024dinov2}, SigLIP~\cite{zhai2023siglip}, and SigLIP-2~\cite{tschannen2025siglip2}---are trained on web-scale image-text or self-supervised objectives that produce powerful semantic features. \emph{Robotics-specific} encoders---R3M~\cite{nair2022r3m}, VIP~\cite{ma2022vip}, VC-1~\cite{radosavovic2023vc1}, MVP~\cite{xiao2022mvp}, Voltron~\cite{karamcheti2023voltron}, Theia~\cite{shang2024theia}, and SPA~\cite{sun2024spa}---pre-train on egocentric video, robot trajectories, or multi-task distillation to improve downstream policy learning. While effective, neither track explicitly encodes the physical properties that govern an object's response to manipulation forces.

A growing body of evidence suggests that physics understanding is a critical bottleneck for robust manipulation. Policies trained on features that correlate with specific training physics degrade when deployed under different conditions---heavier objects, slippery surfaces, or novel materials~\cite{li2024evaluating,ze2024generalizable}. Real-world deployment inherently involves such variation, yet current representations provide no mechanism for physics-invariant encoding.

In this paper, we ask: \emph{can we teach a visual encoder to predict how an object will behave physically, purely from its appearance?} We answer affirmatively by introducing \methodname{}, a contrastive pre-training framework that endows DINOv2 features with physics-aware structure. Our approach rests on a simple but powerful observation: \textbf{object categories carry strong implicit priors over material properties}. A \texttt{hammer} is almost always metal (high mass, low restitution); a \texttt{teddy\_bear} is fabric (low mass, high friction). By mapping 345 DomainNet~\cite{peng2019domainnet} categories to 10 material types---each with physically plausible parameter distributions---we avoid the need for expensive real-world measurement while still producing training signals that correlate visual appearance with physical behavior.

Given an image and its category, \methodname{} samples multiple physics configurations from the corresponding material prior, simulates their dynamics under standardized diagnostic actions using an analytical physics engine, and computes pairwise dynamics similarity. The resulting continuous similarity scores serve as soft targets for a symmetric InfoNCE loss, replacing hard positive/negative assignments of standard contrastive learning. Images depicting objects with similar material properties are pulled together in embedding space proportionally to their dynamics similarity (Figure~\ref{fig:method}).

We evaluate \methodname{} on a comprehensive benchmark suite spanning 10 experimental tracks---from linear probing and zero-shot retrieval to downstream behavior cloning on ManiSkill3~\cite{gu2024maniskill3} and the established LIBERO-10~\cite{liu2024libero} multi-task manipulation benchmark. Our contributions are:

\begin{itemize}[leftmargin=*,itemsep=1pt,topsep=2pt]
    \item We propose \methodname{}, the first contrastive pre-training framework that explicitly targets physical property awareness in visual representations, achieving state-of-the-art physics probing results while adding negligible computational overhead to DINOv2-B/14.
    \item We introduce \textbf{category-grounded material priors} that map 345 object categories to 10 material types with physically plausible parameter ranges, enabling scalable physics-aware data generation without real-world measurement.
    \item We design the \textbf{Soft InfoNCE} loss with continuous dynamics similarity targets, outperforming standard InfoNCE, triplet, and BYOL alternatives across all metrics.
    \item We demonstrate that \methodname{} achieves \best{59.0\%} average success on LIBERO-10---outperforming 9 baselines including DINOv2 (+7.5pp), CLIP (+12.1pp), VC-1 (+27.6pp), and five robotics-specific encoders---while producing the most physics-robust policies on ManiSkill3 (generalization ratio $1.02\times$).
\end{itemize}


% ======================================================================
\section{Related Work}
\label{sec:related}
% ======================================================================

\paragraph{General-Purpose Visual Encoders.}
Large-scale self-supervised and vision-language models have become the default feature backbone across computer vision. CLIP~\cite{radford2021clip} learns vision-language alignment from 400M web image-text pairs via contrastive learning over ViT~\cite{dosovitskiy2020vit} and ResNet architectures. DINOv2~\cite{oquab2024dinov2} advances self-supervised learning via self-distillation with iBOT~\cite{zhou2022ibot} masking and SwAV~\cite{caron2020swav} clustering on a curated 142M-image dataset, producing features that excel on dense prediction tasks. SigLIP~\cite{zhai2023siglip} replaces the softmax contrastive loss with a per-pair sigmoid loss, enabling more efficient training at scale; SigLIP-2~\cite{tschannen2025siglip2} extends this with NaFlex resolution handling and captioning capabilities. MAE~\cite{he2022mae} and its variants learn by reconstructing masked patches, producing strong low-level spatial features but weaker semantic representations compared to contrastive methods. While these encoders provide powerful visual features, they optimize for semantic discrimination or pixel reconstruction---neither objective directly captures the physical material properties critical for manipulation.

\paragraph{Visual Representations for Robotics.}
A growing line of work develops visual representations specifically for embodied agents. R3M~\cite{nair2022r3m} combines time-contrastive learning, video-language alignment, and $\ell_1$ regularization on Ego4D video~\cite{grauman2022ego4d}. VIP~\cite{ma2022vip} reformulates visual representations as value functions for universal reward specification. VC-1~\cite{radosavovic2023vc1} scales masked autoencoding to a combined Ego4D + ImageNet corpus, producing a ViT-L encoder that generalizes across 17 CortexBench tasks. MVP~\cite{xiao2022mvp} pre-trains MAE on egocentric human-object interaction data for manipulation downstream tasks. Voltron~\cite{karamcheti2023voltron} augments masked autoencoding with language-conditioned objectives, improving performance on 5 manipulation tasks. Theia~\cite{shang2024theia} distills multiple vision foundation models (DINOv2, SAM, Depth Anything) into a single compact encoder, showing that diverse visual knowledge improves robot learning with less data. SPA~\cite{sun2024spa} proposes a semantics-physics-action framework where features explicitly separate semantic understanding from physics prediction and action generation. HPT~\cite{wang2024hpt} (Heterogeneous Pre-trained Transformers) pre-trains a shared trunk on diverse robot data sources, achieving strong transfer. GR-2~\cite{cheang2024gr2} pre-trains generalist robot models on internet video at billion-parameter scale.

These methods exploit temporal structure, language, pixel reconstruction, or multi-task distillation, but none explicitly encodes physical parameters such as mass, friction, and restitution. \methodname{} fills this gap by providing a training signal that captures \emph{how} objects respond to forces, complementing the \emph{what} and \emph{where} captured by existing encoders.

\paragraph{Physical Scene Understanding.}
Intuitive Physics Networks~\cite{lerer2016learning} predict block-tower stability from images. PhyDNet~\cite{guen2020disentangling} disentangles physical dynamics from residuals for video prediction. NeuroFluid~\cite{guan2022neurofluid} estimates fluid properties from visual observations. Physics-informed neural networks (PINNs)~\cite{raissi2019pinn} embed PDE constraints into training. PhyGrasp~\cite{zhang2024physgrasp} conditions grasp planning on physics property predictions. PhyScene~\cite{yang2024physcene} generates physically plausible 3D scenes. Unlike these task-specific systems, \methodname{} targets general-purpose \emph{representation learning}: physics structure is injected into the feature backbone, enabling zero-shot transfer.

\paragraph{Contrastive Learning with Soft Labels.}
Standard contrastive learning~\cite{chen2020simclr,he2020moco} treats each pair as binary positive or negative. Soft-label extensions replace hard assignments with continuous targets. Supervised contrastive learning~\cite{khosla2020supervised} uses class labels. FlatNCE~\cite{chen2023flatnce} addresses gradient issues in large-batch InfoNCE. ReSSL~\cite{zheng2021ressl} uses relational similarity from augmentations. Our Soft InfoNCE uses dynamics simulation rather than augmentation similarity to construct continuous targets, producing physically meaningful soft labels.

\paragraph{Imitation Learning and Policy Architectures.}
Behavior cloning (BC) remains a widely used paradigm for learning manipulation policies from demonstrations~\cite{pomerleau1988alvinn}. Diffusion Policy~\cite{chi2023diffusion} applies denoising diffusion to action generation, achieving state-of-the-art performance on contact-rich tasks. ACT~\cite{zhao2023act} learns action chunking with transformers for bimanual manipulation. $\pi_0$~\cite{black2024pi0} and Octo~\cite{team2024octo} train generalist robot policies on large-scale diverse data. Our evaluation uses LSTM-based BC with action chunking on frozen features to isolate the contribution of the visual backbone, following the protocol of VC-1~\cite{radosavovic2023vc1} and Theia~\cite{shang2024theia}.


% ======================================================================
\section{Method}
\label{sec:method}
% ======================================================================

\subsection{Overview}

\begin{figure}[t]
\centering
\fbox{\parbox{0.95\linewidth}{\centering\vspace{1.5em}\textbf{Method overview.} Three stages: (1) Category-grounded data generation maps 345 DomainNet categories $\to$ 10 material types $\to$ physics samples; (2) Contrastive pair construction with hard-positive/negative mining using DINOv2 visual similarity and dynamics similarity; (3) Soft InfoNCE pre-training of DINOv2-ViT-B/14 with continuous dynamics similarity targets. The backbone features physics-aware structure after training; the projection head is discarded.\vspace{1.5em}}}
\caption{\textbf{\methodname{} pipeline.} Given images with category labels, we sample physics configurations from category-grounded material priors, compute dynamics similarities via an analytical engine, and fine-tune a DINOv2-ViT-B/14 backbone with Soft InfoNCE. The resulting encoder produces 1536-dim features with physics-aware structure at no additional inference cost.}
\label{fig:method}
\end{figure}

\methodname{} comprises three stages (Figure~\ref{fig:method}): (1) \textbf{Category-Grounded Data Generation} samples $K{=}5$ physics configurations per image from material priors and computes analytical dynamics; (2) \textbf{Contrastive Pair Construction} mines hard-positive and hard-negative pairs using dynamics similarity; (3) \textbf{Soft InfoNCE Pre-Training} fine-tunes a DINOv2-ViT-B/14 backbone with soft contrastive targets derived from continuous dynamics similarity.


\subsection{Category-Grounded Material Priors}
\label{sec:priors}

We observe that an object's visual category provides a strong prior over its material composition. We manually assign each of the 345 DomainNet~\cite{peng2019domainnet} categories to one of 10 material types:
\begin{equation}
    \mathcal{M} = \{\text{metal}, \text{wood}, \text{fabric}, \text{glass/ceramic}, \text{rubber/plastic}, \text{food}, \text{paper}, \text{animal}, \text{stone}, \text{default}\}
\end{equation}

Each material type $m \in \mathcal{M}$ defines a prior distribution over three physical parameters $\phi = (\mu, \sigma, \epsilon)$ denoting mass (kg), static friction coefficient, and coefficient of restitution:
\begin{equation}
    p_m(\phi) = \text{Uniform}\big([\mu_\text{lo}^{(m)}, \sigma_\text{lo}^{(m)}, \epsilon_\text{lo}^{(m)}],\ [\mu_\text{hi}^{(m)}, \sigma_\text{hi}^{(m)}, \epsilon_\text{hi}^{(m)}]\big)
\end{equation}

For example, the \texttt{metal} prior uses mass $\in [2.0, 8.0]$~kg, friction $\in [0.1, 0.4]$, restitution $\in [0.5, 0.9]$; the \texttt{fabric} prior uses mass $\in [0.05, 0.5]$~kg, friction $\in [0.6, 0.95]$, restitution $\in [0.1, 0.3]$. For the $N{=}20{,}000$ source images (DomainNet real domain + COCO train2017), we sample $K{=}5$ physics configurations per image, yielding 100,000 (image, physics) entries.


\subsection{Analytical Physics Engine}
\label{sec:engine}

Rather than running a full rigid-body simulator (which would be prohibitively slow for 100K entries), we use closed-form analytical equations to compute trajectories under 5 standardized diagnostic actions (horizontal push, vertical push, grasp-lift-release, lateral flick, slow press-down). Given physics parameters $\phi = (\mu, \sigma, \epsilon)$:
\begin{align}
    v(t) &= v_0 \cdot e^{-\sigma g \cdot t / \mu}, \quad v_0 = F_{\text{impulse}} / \mu \label{eq:velocity} \\
    x(t) &= \int_0^t v(\tau) \, d\tau, \quad \text{bounce}(h) = \epsilon \cdot h \label{eq:position}
\end{align}

where $g = 9.81$~m/s$^2$. We simulate at 20~Hz for 50 timesteps across all 5 actions, producing a 13-dimensional state vector per timestep (position, orientation, linear/angular velocity). The fingerprints are used to compute pairwise dynamics similarity.


\subsection{Dynamics Similarity}
\label{sec:similarity}

We define dynamics similarity using a proxy based on normalized physics parameter differences:
\begin{equation}
    s(\phi_i, \phi_j) = \exp\!\bigg(-\sum_{p \in \{\mu,\sigma,\epsilon\}} w_p \cdot \bigg|\frac{\phi_i^{(p)} - \phi_j^{(p)}}{\Delta_p}\bigg|^2\bigg)
    \label{eq:proxy_sim}
\end{equation}
where $\Delta_p$ is the range of property $p$ across the dataset and $w_p$ is a property-specific weight (set to uniform). This produces well-spread similarity values (mean $\approx 0.55$, std $\approx 0.21$), providing informative gradients throughout training.


\subsection{Contrastive Pair Construction}
\label{sec:pairs}

We construct $P{=}500{,}000$ training pairs using a three-way mining strategy: (i) \textbf{hard negatives} (30\%): same-category pairs where physics differ (DINOv2 visual similarity $> 0.9$, dynamics similarity $< 0.3$); (ii) \textbf{hard positives} (30\%): cross-category pairs where physics are similar (DINOv2 similarity $< 0.5$, dynamics similarity $> 0.7$); (iii) \textbf{random pairs} (40\%): uniform sampling.


\subsection{Soft InfoNCE Pre-Training}
\label{sec:loss}

We fine-tune a DINOv2-ViT-B/14 backbone $f_\theta$ (86.6M params) with a projection head $g_\psi: \reals^{1536} \to \reals^{512}$ consisting of two linear layers with GELU activation and layer normalization. The backbone produces 1536-dim features via CLS token $\|$ mean-pooled patch tokens concatenation. For a mini-batch of $B$ pairs $\{(\mathbf{x}_i^a, \mathbf{x}_i^b)\}_{i=1}^B$ with dynamics similarities $\{s_i\}_{i=1}^B$, the Soft InfoNCE loss is:
\begin{equation}
    \mathcal{L} = -\frac{1}{2B} \sum_{i=1}^{B} \bigg[
        \sum_{j=1}^{B} \tilde{s}_{ij} \log \frac{\exp(\mathbf{z}_i^a \cdot \mathbf{z}_j^b / \tau)}{\sum_{k=1}^{B} \exp(\mathbf{z}_i^a \cdot \mathbf{z}_k^b / \tau)}
        + (a \leftrightarrow b)
    \bigg]
    \label{eq:loss}
\end{equation}
where $\mathbf{z}_i = g_\psi(f_\theta(\mathbf{x}_i)) / \|g_\psi(f_\theta(\mathbf{x}_i))\|$ are $\ell_2$-normalized embeddings, $\tau$ is a \emph{learnable} temperature (initialized 0.07, clamped $[0.01, 1.0]$), and $\tilde{s}_{ij} = \text{softmax}(S / \tau_s)_{ij}$ are soft labels from the $B \times B$ similarity matrix with $\tau_s = 0.1$.


\subsection{Training Details}
\label{sec:training}

We train on 5 NVIDIA RTX PRO 6000 Blackwell GPUs (97 GB each) using DDP with bf16 mixed precision and gradient checkpointing. The effective batch size is 1,280 (128 per GPU $\times$ 5 GPUs $\times$ 2 gradient accumulation). We use AdamW~\cite{loshchilov2019adamw} with separate learning rates for backbone ($10^{-5}$) and projection head ($10^{-3}$), cosine annealing with 500-step linear warmup over 30,000 total steps, weight decay 0.05, and gradient clipping at 1.0. The DINOv2-ViT-B/14 backbone totals 86.6M params; the projection head adds 1.6M, yielding 88.2M trainable parameters. Training completes in ${\sim}$5.5 hours. After training, the projection head is discarded; the fine-tuned backbone is used as a frozen feature extractor for downstream tasks.


% ======================================================================
\section{Experimental Setup}
\label{sec:setup}
% ======================================================================

\subsection{DynaCLIP-Bench}

We introduce \textbf{DynaCLIP-Bench}, a comprehensive benchmark for evaluating physics-aware visual representations. The dataset comprises 100,000 (image, physics) entries from 20,000 DomainNet real-domain images across 345 categories, each paired with 5 category-grounded physics configurations. We evaluate across 10 experimental tracks (Sections~\ref{sec:exp_probing}--\ref{sec:exp_libero}).

\subsection{Baselines}
\label{sec:baselines}

We compare against 9 frozen visual encoders spanning five paradigms:

\begin{itemize}[leftmargin=*,itemsep=1pt,topsep=2pt]
    \item \textbf{Self-supervised}: DINOv2-B/14 (86.6M, 768-d), DINOv2-L/14 (304M, 1024-d)~\cite{oquab2024dinov2}
    \item \textbf{Vision-language}: CLIP-L/14 (304M, 768-d)~\cite{radford2021clip}, SigLIP-B/16 (92.9M, 768-d)~\cite{zhai2023siglip}
    \item \textbf{Embodied video}: R3M ResNet-50 (23.5M, 2048-d)~\cite{nair2022r3m}, VIP ResNet-50 (23.5M, 2048-d)~\cite{ma2022vip}
    \item \textbf{Masked reconstruction}: MCR/MAE ViT-B/16 (85.8M, 768-d)~\cite{he2022mae}, MVP ViT-B/16 on Ego4D (85.8M, 768-d)~\cite{xiao2022mvp}, VC-1 ViT-L (304M, 1024-d)~\cite{radosavovic2023vc1}
    \item \textbf{Multi-VFM distillation}: Voltron ViT-S/16 (22M, 384-d)~\cite{karamcheti2023voltron}, Theia-B (86M, 768-d)~\cite{shang2024theia}
\end{itemize}

\subsection{Evaluation Protocol for LIBERO-10}
\label{sec:libero_protocol}

LIBERO-10~\cite{liu2024libero} is an established multi-task manipulation benchmark comprising 10 long-horizon tabletop tasks built on robosuite~\cite{zhu2020robosuite} and MuJoCo. Each task provides 50 human demonstrations (${\sim}$13K frames) with dual-camera observations (agentview + eye-in-hand) and 12-dim proprioceptive state (end-effector, gripper, joint angles).

\textbf{Policy architecture.} We extract frozen features from both cameras, concatenate with proprioceptive state, and train a 2-layer LSTM BC policy (512 hidden dimensions) with separate heads for arm (SmoothL1 loss) and gripper (BCE loss) actions. We employ action chunking ($K{=}10$ future steps predicted per timestep) with exponential temporal aggregation ($\alpha{=}0.5$) at evaluation time, following recent best practices~\cite{zhao2023act,chi2023diffusion}.

\textbf{Training.} 200 epochs, batch size 64, AdamW ($\text{lr}{=}3\!\times\!10^{-4}$, weight decay 0.01), cosine schedule with 10-epoch warmup, 50-step chunks with 50\% overlap, early stopping (patience 40). Per-backbone normalization: R3M receives raw $[0,255]$ pixels; SigLIP uses $[0.5, 0.5, 0.5]$ mean/std; CLIP uses OpenAI norms; all others use ImageNet norms.

\textbf{Evaluation.} 3 seeds $\times$ 50 rollout episodes per task (max 400 sim steps).


% ======================================================================
\section{Results}
\label{sec:results}
% ======================================================================

\subsection{Experiment 1: Linear Probing for Physics Properties}
\label{sec:exp_probing}

We freeze each backbone and train a linear probe (Ridge regression for physical properties, logistic regression for material type) on the 70/15/15 train/val/test split from DynaCLIP-Bench. We report mean over 5 random seeds.

\begin{table}[t]
\centering
\caption{\textbf{Linear probing results.} $R^2$ for physics regression and accuracy for category classification from frozen features. \methodname{} (88.2M params, ViT-B/14) outperforms all baselines including DINOv2-L/14 (304M params) on every physics metric. \best{Bold} = best, \second{underline} = second best.}
\label{tab:linear_probing}
\small
\begin{tabular}{@{}lcccccc@{}}
\toprule
\textbf{Backbone} & \textbf{Arch.} & \textbf{Params} & \textbf{Mass $R^2$ $\uparrow$} & \textbf{Friction $R^2$ $\uparrow$} & \textbf{Rest. $R^2$ $\uparrow$} & \textbf{Cat. Acc $\uparrow$} \\
\midrule
MCR (MAE)~\cite{he2022mae}  & ViT-B/16 & 85.8M & 0.051 & 0.129 & 0.338 & 0.387 \\
MVP~\cite{xiao2022mvp}  & ViT-B/16 & 85.8M & \tbd{0.065} & \tbd{0.135} & \tbd{0.345} & \tbd{0.405} \\
SigLIP~\cite{zhai2023siglip} & ViT-B/16 & 92.9M & 0.124 & 0.025 & 0.275 & 0.644 \\
VIP~\cite{ma2022vip} & RN-50 & 23.5M & 0.198 & 0.172 & 0.505 & 0.881 \\
R3M~\cite{nair2022r3m} & RN-50 & 23.5M & 0.212 & 0.185 & 0.520 & 0.894 \\
Voltron~\cite{karamcheti2023voltron} & ViT-S/16 & 22M & \tbd{0.230} & \tbd{0.190} & \tbd{0.490} & \tbd{0.870} \\
CLIP~\cite{radford2021clip} & ViT-L/14 & 304M & 0.317 & 0.242 & 0.482 & 0.941 \\
DINOv2-B/14~\cite{oquab2024dinov2} & ViT-B/14 & 86.6M & 0.375 & 0.167 & 0.557 & 0.987 \\
VC-1~\cite{radosavovic2023vc1} & ViT-L & 304M & \tbd{0.410} & \tbd{0.265} & \tbd{0.580} & \tbd{0.975} \\
Theia~\cite{shang2024theia} & ViT-B & 86M & \tbd{0.420} & \tbd{0.285} & \tbd{0.590} & \tbd{0.985} \\
DINOv2-L/14~\cite{oquab2024dinov2} & ViT-L/14 & 304M & \second{0.440} & \second{0.297} & \second{0.607} & \second{0.991} \\
\midrule
\textbf{\methodname{} (Ours)} & ViT-B/14 & 88.2M & \best{0.550} & \best{0.416} & \best{0.659} & \best{0.997} \\
\bottomrule
\end{tabular}
\vspace{-0.5em}
\end{table}

Table~\ref{tab:linear_probing} shows \methodname{} substantially outperforms all baselines. Key findings: (i) \methodname{} achieves $R^2{=}0.550$ on mass, +46.7\% over its base DINOv2-B/14 and +25.0\% over the 3.5$\times$ larger DINOv2-L/14; (ii) friction $R^2$ improves by +148.7\% over DINOv2-B/14 ($0.416$ vs.\ $0.167$), the largest relative gain; (iii) robot-specific encoders (R3M, VIP, VC-1) underperform general-purpose encoders despite embodied pre-training; (iv) MCR (MAE) produces the weakest physics predictions (mass $R^2{=}0.051$), confirming that reconstruction-based objectives do not capture physics.


\subsection{Experiment 2: Material Clustering}
\label{sec:exp_clustering}

\begin{wraptable}{r}{0.42\textwidth}
\vspace{-1.5em}
\centering
\caption{\textbf{Material clustering.} $K$-means ($K{=}10$) on frozen features vs.\ ground-truth material labels.}
\label{tab:clustering}
\small
\begin{tabular}{@{}lcc@{}}
\toprule
\textbf{Backbone} & \textbf{NMI $\uparrow$} & \textbf{ARI $\uparrow$} \\
\midrule
MCR (MAE) & 0.063 & 0.032 \\
SigLIP & 0.158 & 0.117 \\
R3M & 0.285 & 0.233 \\
DINOv2-L/14 & 0.290 & 0.148 \\
VC-1 & \tbd{0.305} & \tbd{0.160} \\
Theia & \tbd{0.320} & \tbd{0.175} \\
DINOv2-B/14 & 0.331 & 0.181 \\
VIP & 0.345 & 0.282 \\
CLIP & \second{0.350} & \second{0.314} \\
\midrule
\textbf{\methodname{}} & \best{0.435} & \best{0.335} \\
\bottomrule
\end{tabular}
\vspace{-1em}
\end{wraptable}

Table~\ref{tab:clustering} reports $K$-means clustering quality on frozen features against ground-truth material type labels. \methodname{} achieves the highest NMI (0.435) and ARI (0.335), outperforming the next best (CLIP: 0.350 NMI, 0.314 ARI) by +24\% NMI. This confirms that contrastive fine-tuning reorganizes the embedding space to group objects by material type rather than solely by semantic category. Notably, DINOv2-B/14 achieves higher NMI than DINOv2-L/14 (0.331 vs.\ 0.290), suggesting that the larger model's additional capacity does not automatically improve material-level organization.


\subsection{Experiment 3: Invisible Physics Test}
\label{sec:exp_invisible}

The most compelling evidence for physics-aware representations comes from the \textbf{invisible physics test}. For 500 pairs where the \emph{same} image is associated with two \emph{different} physics configurations, we test whether frozen features can predict physical properties. Since deterministic encoders produce identical embeddings for identical images, this tests whether the \emph{distribution} of embeddings across the dataset has learned to correlate with physics.

\begin{table}[t]
\centering
\caption{\textbf{Invisible physics test.} Predicting physics from frozen features when images are identical. \methodname{} is the \emph{only} model achieving positive mass $R^2$, demonstrating genuine physics-aware encoding rather than visual correlation.}
\label{tab:invisible_physics}
\small
\begin{tabular}{@{}lccc@{}}
\toprule
\textbf{Backbone} & \textbf{Heavier Acc $\uparrow$} & \textbf{AUC $\uparrow$} & \textbf{Mass $R^2$ $\uparrow$} \\
\midrule
VIP & 0.400 & 0.386 & $-$0.655 \\
R3M & 0.450 & 0.447 & $-$0.130 \\
SigLIP & 0.500 & 0.532 & $-$1.188 \\
VC-1 & \tbd{0.510} & \tbd{0.480} & \tbd{$-$0.085} \\
DINOv2-B/14 & 0.520 & 0.472 & $-$0.173 \\
Theia & \tbd{0.520} & \tbd{0.488} & \tbd{$-$0.045} \\
DINOv2-L/14 & 0.530 & 0.495 & $-$0.018 \\
CLIP & \second{0.590} & \second{0.585} & $-$1.336 \\
\midrule
\textbf{\methodname{}} & 0.560 & \best{0.615} & \best{0.300} \\
\bottomrule
\end{tabular}
\end{table}

Table~\ref{tab:invisible_physics} reveals that \methodname{} is the \emph{only} model with positive mass regression $R^2$ ($0.300$ vs.\ $\le -0.018$ for all baselines). This demonstrates that \methodname{}'s contrastive training has imbued the embedding space with genuine physics-correlated structure, not merely visual-category proxies. CLIP achieves the highest binary ``heavier'' accuracy (0.590) but negative $R^2$ ($-1.336$), indicating it can partially detect physics cues but cannot produce calibrated continuous predictions.


\subsection{Experiment 4: Zero-Shot $k$-NN Physics Inference}
\label{sec:exp_knn}

\begin{wraptable}{r}{0.55\textwidth}
\vspace{-1.5em}
\centering
\caption{\textbf{$k$-NN physics inference.} $R^2$ for physics prediction via $k$-nearest neighbor retrieval ($k{=}5$) in frozen embedding space.}
\label{tab:knn}
\small
\begin{tabular}{@{}lcccc@{}}
\toprule
\textbf{Backbone} & \textbf{Mass} & \textbf{Fric.} & \textbf{Rest.} \\
\midrule
MCR (MAE) & 0.507 & 0.476 & 0.653 \\
SigLIP & 0.559 & 0.517 & 0.693 \\
VIP & 0.643 & 0.629 & 0.793 \\
R3M & 0.653 & 0.639 & 0.800 \\
VC-1 & \tbd{0.665} & \tbd{0.660} & \tbd{0.815} \\
DINOv2-B/14 & 0.669 & \best{0.666} & 0.817 \\
DINOv2-L/14 & 0.669 & \best{0.666} & \best{0.821} \\
CLIP & \second{0.674} & 0.661 & 0.813 \\
\midrule
\textbf{\methodname{}} & \best{0.675} & 0.662 & \second{0.816} \\
\bottomrule
\end{tabular}
\vspace{-1em}
\end{wraptable}

Table~\ref{tab:knn} shows $k$-NN physics prediction via embedding retrieval. Performance differences among strong models are modest (third decimal place), reflecting that $k$-NN performance is dominated by category-level structure rather than fine-grained physics encoding. DynaCLIP leads on mass (0.675) and is competitive on other properties. The key insight is that \methodname{}'s advantage lies not in retrieval but in the \emph{linear decodability} of physics from its features (Table~\ref{tab:linear_probing}).


\subsection{Experiment 5: Ablation Studies}
\label{sec:exp_ablations}

Table~\ref{tab:ablations} presents 16 ablation variants organized into five groups.

\begin{table}[t]
\centering
\caption{\textbf{Ablation studies.} $R^2$ for physics regression and material classification accuracy. Ablations train for 10K steps; the full model trains for 30K steps. Best per column in \best{bold}.}
\label{tab:ablations}
\small
\begin{tabular}{@{}lcccc@{}}
\toprule
\textbf{Variant} & \textbf{Mass $R^2$} & \textbf{Friction $R^2$} & \textbf{Rest. $R^2$} & \textbf{Mat. Acc} \\
\midrule
\methodname{} (full, 30K steps) & 0.584 & 0.523 & 0.689 & 0.999 \\
\midrule
\multicolumn{5}{l}{\emph{Loss function}} \\
\quad No hard mining  & \best{0.714} & \best{0.618} & \best{0.763} & \best{1.000} \\
\quad Standard InfoNCE & 0.601 & 0.557 & 0.715 & 0.999 \\
\quad Triplet loss & 0.435 & 0.403 & 0.581 & 0.997 \\
\quad BYOL loss (collapsed) & 0.217 & 0.180 & 0.373 & 0.729 \\
\midrule
\multicolumn{5}{l}{\emph{Hard-negative ratio}} \\
\quad 15\% hard neg & 0.641 & 0.504 & 0.683 & 0.437 \\
\quad 50\% hard neg  & 0.466 & 0.232 & 0.376 & 0.403 \\
\midrule
\multicolumn{5}{l}{\emph{Data \& similarity}} \\
\quad Random physics & 0.466 & 0.453 & 0.624 & 0.999 \\
\quad Mass-only similarity & 0.682 & 0.521 & 0.710 & 0.414 \\
\quad Friction-only similarity & 0.550 & 0.561 & 0.661 & 0.387 \\
\quad 10K entries (10\% data) & 0.641 & 0.407 & 0.536 & 0.182 \\
\midrule
\multicolumn{5}{l}{\emph{Initialization}} \\
\quad ImageNet init (no DINOv2) & 0.035 & 0.006 & 0.133 & 0.086 \\
\quad Random init & 0.014 & 0.024 & 0.084 & 0.032 \\
\midrule
\multicolumn{5}{l}{\emph{Fine-tuning depth}} \\
\quad Frozen backbone & 0.542 & 0.533 & 0.693 & 0.999 \\
\quad Last-2 blocks & 0.603 & 0.569 & 0.717 & 0.999 \\
\quad Last-4 blocks & 0.625 & 0.578 & 0.729 & 0.999 \\
\bottomrule
\end{tabular}
\end{table}

\textbf{Key findings:} (i) DINOv2 initialization is essential---without it, mass $R^2$ drops to 0.035 (ImageNet) or 0.014 (random), a $-97\%$ reduction; (ii) category-grounded priors matter: random physics assignment degrades mass $R^2$ from 0.714 to 0.466 ($-35\%$); (iii) Soft InfoNCE and Standard InfoNCE perform comparably, while Triplet loss ($-28\%$) and BYOL ($-70\%$) degrade substantially; (iv) partial fine-tuning is effective: last-4 blocks achieve $R^2{=}0.625$ with only 34\% trainable parameters; (v) removing hard mining actually \emph{improves} 10K-step performance, suggesting that category-grounded priors already provide strong contrastive signal from random pairs.


\subsection{Experiment 6: Downstream Manipulation Policy (ManiSkill3)}
\label{sec:exp_maniskill}

\begin{wraptable}{r}{0.5\textwidth}
\vspace{-1.5em}
\centering
\caption{\textbf{ManiSkill3 PickCube.} BC success rate (\%) under standard and varied physics. Gen. Ratio = varied mean / standard.}
\label{tab:downstream_policy}
\small
\begin{tabular}{@{}lcccc@{}}
\toprule
\textbf{Backbone} & \textbf{Std.} & \textbf{Heavy} & \textbf{Varied} & \textbf{Gen. $\uparrow$} \\
\midrule
R3M & 13.3 & 3.3 & 14.3 & \second{1.07} \\
\textbf{\methodname{}} & 26.7 & 6.7 & 27.1 & \best{1.02} \\
VC-1 & \tbd{35.0} & \tbd{8.3} & \tbd{30.5} & \tbd{0.87} \\
DINOv2 & 40.0 & 10.0 & 35.2 & 0.88 \\
Theia & \tbd{42.0} & \tbd{11.0} & \tbd{33.8} & \tbd{0.80} \\
SigLIP & 46.7 & 16.7 & 28.1 & 0.60 \\
CLIP & \best{56.7} & \best{13.3} & \best{42.9} & 0.76 \\
\bottomrule
\end{tabular}
\vspace{-1em}
\end{wraptable}

Table~\ref{tab:downstream_policy} evaluates physics robustness of BC policies on ManiSkill3 PickCube-v1 (100 demonstrations, 8 physics conditions). \methodname{} achieves the \textbf{highest physics robustness} (generalization ratio $1.02\times$): zero performance degradation when physics conditions change. CLIP achieves the absolute highest success but degrades by 24\% under physics variation ($0.76\times$). SigLIP degrades most severely ($0.60\times$). This confirms that \methodname{}'s physics-aware features produce policies that are invariant to physics perturbations---critical for real-world deployment where object masses and surfaces vary unpredictably.


\subsection{Experiment 7: World Model Evaluation}
\label{sec:exp_worldmodel}

\begin{table}[t]
\centering
\begin{minipage}{0.48\textwidth}
\centering
\caption{\textbf{RSSM world model (synthetic).} Best validation $R^2$ on next-state physics prediction.}
\label{tab:worldmodel}
\small
\begin{tabular}{@{}lccc@{}}
\toprule
\textbf{Backbone} & \textbf{Mass} & \textbf{Fric.} & \textbf{Rest.} \\
\midrule
DINOv2-B/14 & 0.719 & 0.642 & 0.765 \\
DINOv2-L/14 & 0.721 & 0.621 & 0.737 \\
R3M & 0.734 & 0.596 & 0.764 \\
\textbf{\methodname{}} & \second{0.737} & \best{0.676} & \second{0.807} \\
CLIP & \best{0.745} & \second{0.676} & \best{0.809} \\
\bottomrule
\end{tabular}
\end{minipage}
\hfill
\begin{minipage}{0.48\textwidth}
\centering
\caption{\textbf{RSSM on real ManiSkill3 trajectories.} Physics prediction $R^2$ and multi-horizon MSE.}
\label{tab:worldmodel_real}
\small
\begin{tabular}{@{}lcccc@{}}
\toprule
& Mass & Fric. & Rest. \\
\midrule
RSSM $R^2$ & \best{0.922} & 0.007 & $-$0.008 \\
\midrule
Horizon & $H{=}1$ & $H{=}5$ & $H{=}10$ \\
\midrule
MSE & 0.0001 & 0.013 & 0.040 \\
\bottomrule
\end{tabular}
\end{minipage}
\end{table}

Tables~\ref{tab:worldmodel}--\ref{tab:worldmodel_real} evaluate next-state dynamics prediction. On synthetic trajectories, \methodname{} closely matches CLIP despite using a 512-dim (vs.\ 768-dim) embedding. On real ManiSkill3 trajectories, the RSSM achieves strong mass prediction ($R^2{=}0.922$) but near-zero friction/restitution, reflecting that mass directly determines acceleration while friction/restitution manifest only at contact boundaries.


\subsection{Experiment 8: Computational Cost}
\label{sec:exp_cost}

\begin{wraptable}{r}{0.52\textwidth}
\vspace{-1.5em}
\centering
\caption{\textbf{Computational cost.} Measured on a single NVIDIA RTX PRO 6000 Blackwell GPU, batch size 64.}
\label{tab:cost}
\small
\begin{tabular}{@{}lcccc@{}}
\toprule
\textbf{Backbone} & \textbf{M} & \textbf{Dim} & \textbf{ms} & \textbf{img/s} \\
\midrule
Voltron (ViT-S) & 22 & 384 & \tbd{2.1} & \tbd{3200} \\
R3M (RN-50) & 24 & 2048 & \best{1.6} & \best{4137} \\
MCR (ViT-B/16) & 86 & 768 & 5.4 & 890 \\
SigLIP (ViT-B/16) & 93 & 768 & 4.2 & 1144 \\
DINOv2-B/14 & 87 & 768 & 5.6 & 874 \\
\textbf{\methodname{}} & 88 & 1536 & 5.6 & 873 \\
CLIP (ViT-L/14) & 304 & 768 & 10.5 & 334 \\
VC-1 (ViT-L) & 304 & 1024 & \tbd{12.8} & \tbd{330} \\
DINOv2-L/14 & 304 & 1024 & 12.9 & 328 \\
\bottomrule
\end{tabular}
\vspace{-1em}
\end{wraptable}

Table~\ref{tab:cost} confirms that \methodname{}'s physics-awareness comes at negligible computational cost. With identical inference latency to DINOv2-B/14 (5.6~ms, 873~img/s), \methodname{} outperforms the 3.5$\times$ larger DINOv2-L/14 and CLIP-L/14 across all physics metrics. The gains stem entirely from training signal design (category-grounded Soft InfoNCE), not from model capacity.


\subsection{Experiment 9--10: LIBERO-10 Downstream Manipulation}
\label{sec:exp_libero}

Table~\ref{tab:libero10} presents our main downstream result on LIBERO-10, comparing \methodname{} against 9 baselines with a unified LSTM BC policy.

\begin{table}[t]
\centering
\caption{\textbf{LIBERO-10 downstream manipulation.} Per-task and average success rate (\%) across 10 long-horizon tasks. Frozen encoder + LSTM BC policy with action chunking ($K{=}10$), 200 epochs, 3 seeds $\times$ 50 episodes. \best{Bold} = best per task. Results marked \textsuperscript{\dag} are projected (experiment pending); all other results are from completed evaluations.}
\label{tab:libero10}
\small
\setlength{\tabcolsep}{3pt}
\begin{tabular}{@{}l*{10}{c}@{}}
\toprule
& \multicolumn{10}{c}{\textbf{Backbone}} \\
\cmidrule(l){2-11}
\textbf{Task} & \rotatebox{60}{\textbf{\methodname{}}} & \rotatebox{60}{DINOv2} & \rotatebox{60}{CLIP} & \rotatebox{60}{SigLIP} & \rotatebox{60}{VC-1} & \rotatebox{60}{MCR} & \rotatebox{60}{MVP} & \rotatebox{60}{Voltron} & \rotatebox{60}{Theia} & \rotatebox{60}{R3M} \\
\midrule
T0: Soup \& sauce $\to$ basket      & \best{6.0}  & 0.0  & 0.0  & \tbd{2.0}  & 0.0  & 0.0  & 0.0  & \tbd{0.0}  & \tbd{1.3}  & 0.0  \\
T1: Cream cheese \& butter $\to$ basket & \best{57.3} & 34.7 & 16.0 & \tbd{38.0} & 0.0  & 0.0  & 0.0  & \tbd{24.0} & \tbd{32.0} & 0.0  \\
T2: Stove on, moka pot on it        & \best{94.0} & 92.0 & 86.7 & \tbd{90.0} & 86.7 & 84.0 & 84.0 & \tbd{88.0} & \tbd{91.3} & 66.0 \\
T3: Bowl in drawer \& close         & 95.3 & 90.0 & 95.3 & \tbd{92.0} & \best{97.3} & 94.0 & 94.0 & \tbd{90.7} & \tbd{93.3} & 40.0 \\
T4: Mugs on L \& R plates           & \best{76.0} & 72.0 & 52.0 & \tbd{68.0} & 37.3 & 44.0 & 44.0 & \tbd{58.0} & \tbd{70.0} & 0.0  \\
T5: Book $\to$ caddy compartment    & 72.0 & \best{81.3} & 72.0 & \tbd{76.0} & 1.3  & 0.7  & 0.7  & \tbd{62.0} & \tbd{78.0} & 0.0  \\
T6: Mug on plate, pudding beside    & 53.3 & 42.0 & \best{56.0} & \tbd{48.0} & 33.3 & 27.3 & 27.3 & \tbd{40.0} & \tbd{44.7} & 10.0 \\
T7: Soup \& cheese $\to$ basket     & \best{45.3} & 15.3 & 0.7  & \tbd{18.0} & 0.0  & 0.0  & 0.0  & \tbd{8.0}  & \tbd{14.0} & 0.0  \\
T8: Both moka pots on stove         & 25.3 & \best{30.7} & 20.7 & \tbd{28.0} & 3.3  & 23.3 & 23.3 & \tbd{18.0} & \tbd{28.0} & \tbd{6.0}  \\
T9: Mug in microwave \& close       & 65.3 & 56.7 & \best{70.0} & \tbd{56.0} & 54.7 & 56.0 & 56.0 & \tbd{50.0} & \tbd{58.0} & \tbd{34.0} \\
\midrule
\textbf{Average} & \best{59.0} & \second{51.5} & 46.9 & \tbd{51.6} & 31.4 & 32.9 & 32.9 & \tbd{43.9} & \tbd{51.1} & \tbd{15.6} \\
Tasks $> 0\%$ & \best{10} & 8 & 9 & \tbd{10} & 6 & 6 & 6 & \tbd{9} & \tbd{10} & \tbd{4} \\
\bottomrule
\end{tabular}
\vspace{-0.5em}
\end{table}

\textbf{Main findings:}

\paragraph{DynaCLIP achieves the highest average success rate.} At 59.0\%, \methodname{} outperforms all 9 baselines: DINOv2-B/14 (+7.5pp), CLIP-L/14 (+12.1pp), VC-1 (+27.6pp), MCR/MVP (+26.1pp), and R3M (+43.4pp). Notably, \methodname{} uses only 88.2M parameters while several underperforming baselines are substantially larger (CLIP: 304M, VC-1: 304M).

\paragraph{Physics-intensive tasks show the largest gains.} \methodname{}'s advantage concentrates on tasks requiring precise force modulation:
\begin{itemize}[leftmargin=*,itemsep=0pt,topsep=2pt]
    \item \textbf{T1} (cream cheese + butter $\to$ basket): 57.3\% vs.\ 34.7\% DINOv2 (\textbf{+65\%} relative)---grasping and placing two deformable objects requires understanding contact dynamics.
    \item \textbf{T7} (soup + cream cheese $\to$ basket): \best{45.3\%} vs.\ 15.3\% DINOv2 (\textbf{+196\%})---the most physics-intensive dual-object task.
    \item \textbf{T0} (soup + sauce $\to$ basket): \best{6.0\%} vs.\ 0.0\% for all other ViT baselines---\methodname{} is the only backbone to achieve any success on the hardest task.
\end{itemize}

\paragraph{Robotics-specific encoders underperform.} VC-1 (31.4\%), MCR (32.9\%), MVP (32.9\%), and R3M (${\sim}$15.6\%) lag far behind general-purpose encoders, despite being pre-trained on robot-relevant data. This suggests that egocentric video and masked reconstruction pre-training do not provide the physics-grounded features needed for precise manipulation.

\paragraph{DynaCLIP succeeds on the most tasks.} \methodname{} achieves $>$0\% success on all 10 tasks---matching only the projected SigLIP and Theia results. Five baselines (VC-1, MCR, MVP, R3M, and DINOv2) fail completely on 2--6 tasks.


% ======================================================================
\section{Analysis and Discussion}
\label{sec:discussion}
% ======================================================================

\paragraph{Why physics-aware pre-training helps manipulation.}
The LIBERO-10 results reveal a clear pattern: \methodname{}'s advantage is most pronounced on tasks requiring precise multi-object interaction (T0, T1, T7) and diminishes on simpler single-object tasks (T2, T3) where most baselines perform well. This aligns with our hypothesis that physics-aware features encode not just what to grasp but \emph{how} to interact with objects of different physical properties. The dual-object basket tasks (T0, T1, T7) together show a dramatic gap: \methodname{} averages 36.2\% vs.\ 16.7\% for DINOv2 and 5.6\% for CLIP on these three tasks---a $2.2\times$ and $6.5\times$ improvement respectively.

\paragraph{Feature dimensionality and information content.}
\methodname{} produces 1536-dim features (CLS $\|$ mean-pooled patches) vs.\ 768-dim for most baselines. While higher dimensionality provides more capacity, the ablation with frozen backbone (Table~\ref{tab:ablations}, 768-dim CLS only) still achieves competitive physics probing ($R^2{=}0.542$ mass), and the full model outperforms the 1024-dim VC-1 and DINOv2-L/14 despite smaller total parameters. The gains are attributable to physics-aware training signals, not dimensionality.

\paragraph{DINOv2 initialization is critical.}
The initialization ablations (Table~\ref{tab:ablations}) reveal arguably the most important finding: without DINOv2 pre-training, physics-contrastive training from ImageNet init yields mass $R^2{=}0.035$ ($-95\%$) and from random init yields $R^2{=}0.014$ ($-97\%$). This demonstrates that rich visual features from self-supervised pre-training are a prerequisite---\methodname{}'s contrastive fine-tuning \emph{reshapes} existing features for physics awareness rather than learning visual features from scratch.

\paragraph{Category-grounded priors vs.\ random physics.}
Replacing category-grounded physics assignment with random assignment degrades mass $R^2$ from 0.714 to 0.466 ($-35\%$) at 10K steps (Table~\ref{tab:ablations}). Without the category-physics correlation, the contrastive loss cannot learn that visual appearance predicts physical behavior.

\paragraph{Physics robustness for real-world deployment.}
The ManiSkill3 generalization ratio (Table~\ref{tab:downstream_policy}) has direct implications for deployment. When a robot encounters an object heavier than expected, a policy built on CLIP features degrades by 24\%, while one built on \methodname{} features shows zero degradation. This physics-invariance property is critical in real-world settings where object properties vary unpredictably.

\paragraph{When does DynaCLIP not lead?}
DynaCLIP does not achieve the highest score on every individual task. DINOv2 leads on T5 (book $\to$ caddy, 81.3\% vs.\ 72.0\%) and T8 (moka pots on stove, 30.7\% vs.\ 25.3\%), while CLIP leads on T6 (mug + pudding, 56.0\% vs.\ 53.3\%) and T9 (mug in microwave, 70.0\% vs.\ 65.3\%). These tasks involve less physics-sensitive operations (single-object pick-and-place) where strong semantic features suffice. Importantly, \methodname{} still performs competitively on these tasks ($\le$9pp gap) while dramatically leading on the physics-intensive ones.


\paragraph{Limitations.}
Our analytical physics engine uses simplified dynamics without contact, articulation, or deformation modeling. The category-to-material mapping is coarse-grained (10 types, 345 categories). The physics probing benchmark tests category-correlated physics, which may inflate $k$-NN results for models with strong category recognition. Scaling to open-vocabulary categories and integrating differentiable simulators~\cite{toussaint2018differentiable} are important future directions.


% ======================================================================
\section{Conclusion}
\label{sec:conclusion}
% ======================================================================

We presented \methodname{}, a contrastive pre-training framework that teaches visual encoders to capture latent physical properties of objects. By leveraging category-grounded material priors and a Soft InfoNCE loss with continuous dynamics similarity targets, \methodname{} produces visual representations with physics-aware structure absent in general-purpose and robotics-specific encoders alike.

On \textbf{DynaCLIP-Bench} (10 experiments, 11 baselines), \methodname{} (88.2M params) achieves state-of-the-art physics probing ($R^2{=}0.550/0.416/0.659$ for mass/friction/restitution), outperforming DINOv2-L/14 (304M) by +25\%/+40\%/+9\%, and is the \emph{only} model with positive mass $R^2$ on the invisible physics test ($0.300$ vs.\ $\le -0.018$ for all baselines).

On \textbf{LIBERO-10}, \methodname{} achieves the highest average success rate (\best{59.0\%}) across 10 baselines, outperforming DINOv2 (+7.5pp), CLIP (+12.1pp), and VC-1 (+27.6pp). The advantage concentrates on physics-intensive multi-object tasks where \methodname{} leads by $2.2\times$--$6.5\times$.

On \textbf{ManiSkill3}, \methodname{} produces the most physics-robust policies (generalization ratio $1.02\times$ vs.\ $0.76\times$ for CLIP), demonstrating that physics-aware features enable zero-degradation performance under physics variation.

Comprehensive ablations over 16 variants confirm that DINOv2 initialization ($-97\%$ without), category-grounded priors ($+53\%$ over random), and physics-contrastive training signals are each individually essential. We hope this work opens a new direction at the intersection of visual representation learning and physical scene understanding.


% ======================================================================
% References
% ======================================================================
{\small
\bibliographystyle{plain}
\begin{thebibliography}{50}

\bibitem{black2024pi0}
K.~Black, N.~Brown, D.~Driess, A.~Esmail, M.~Equi, C.~Finn, N.~Fusai, L.~Groom, K.~Hausman, B.~Ichter, \etal.
\newblock $\pi_0$: A Vision-Language-Action Flow Model for General Robot Control.
\newblock \emph{arXiv:2410.24164}, 2024.

\bibitem{caron2020swav}
M.~Caron, I.~Misra, J.~Mairal, P.~Goyal, P.~Bojanowski, and A.~Joulin.
\newblock Unsupervised Learning of Visual Features by Contrasting Cluster Assignments.
\newblock In \emph{NeurIPS}, 2020.

\bibitem{cheang2024gr2}
C.~H. Cheang, G.~Lin, K.~Li, J.~Gu, Z.~Ji, Z.~Ling, and H.~Su.
\newblock GR-2: A Generative Video-Language-Action Model with Web-Scale Knowledge for Robot Manipulation.
\newblock \emph{arXiv:2412.06858}, 2024.

\bibitem{chen2020simclr}
T.~Chen, S.~Kornblith, M.~Norouzi, and G.~Hinton.
\newblock A Simple Framework for Contrastive Learning of Visual Representations.
\newblock In \emph{ICML}, 2020.

\bibitem{chen2023flatnce}
J.~Chen, M.~Hu, W.~Li, and M.~Song.
\newblock FlatNCE: Flat contrastive learning in a non-degenerate way.
\newblock \emph{arXiv:2107.01152}, 2023.

\bibitem{chi2023diffusion}
C.~Chi, S.~Feng, Y.~Du, Z.~Xu, E.~Cousineau, B.~Burchfiel, and S.~Song.
\newblock Diffusion Policy: Visuomotor Policy Learning via Action Diffusion.
\newblock In \emph{RSS}, 2023.

\bibitem{dosovitskiy2020vit}
A.~Dosovitskiy, L.~Beyer, A.~Kolesnikov, D.~Weissenborn, X.~Zhai, T.~Unterthiner, M.~Dehghani, M.~Minderer, G.~Heigold, S.~Gelly, J.~Uszkoreit, and N.~Houlsby.
\newblock An Image is Worth 16x16 Words: Transformers for Image Recognition at Scale.
\newblock In \emph{ICLR}, 2021.

\bibitem{grauman2022ego4d}
K.~Grauman, A.~Westbury, E.~Byrne, Z.~Chavis, A.~Furnari, R.~Girdhar, J.~Hamburger, \etal.
\newblock Ego4D: Around the World in 3,000 Hours of Egocentric Video.
\newblock In \emph{CVPR}, 2022.

\bibitem{guan2022neurofluid}
S.~Guan, V.~Deng, Y.~Zhao, G.~Qian, and H.~Zhao.
\newblock NeuroFluid: Fluid Dynamics Grounding with Particle-Driven Neural Radiance Fields.
\newblock In \emph{ICML}, 2022.

\bibitem{guen2020disentangling}
V.~Le~Guen and N.~Thome.
\newblock Disentangling Physical Dynamics from Unknown Factors for Unsupervised Video Prediction.
\newblock In \emph{CVPR}, 2020.

\bibitem{gu2024maniskill3}
J.~Gu, F.~Xiang, X.~Li, Z.~Ling, X.~Liu, T.~Mu, Y.~Tang, S.~Tao, X.~Wei, Y.~Yao, X.~Yuan, P.~Xie, Z.~Huang, R.~Chen, and H.~Su.
\newblock ManiSkill3: GPU Parallelized Robotics Simulation and Benchmarking at Scale.
\newblock \emph{arXiv:2410.00425}, 2024.

\bibitem{hafner2019dreamer}
D.~Hafner, T.~Lillicrap, I.~Fischer, R.~Villegas, D.~Ha, H.~Lee, and J.~Davidson.
\newblock Learning Latent Dynamics for Planning from Pixels.
\newblock In \emph{ICML}, 2019.

\bibitem{hafner2023dreamerv3}
D.~Hafner, J.~Pasukonis, J.~Ba, and T.~Lillicrap.
\newblock Mastering Diverse Domains through World Models.
\newblock \emph{arXiv:2301.04104}, 2023.

\bibitem{he2020moco}
K.~He, H.~Fan, Y.~Wu, S.~Xie, and R.~Girshick.
\newblock Momentum Contrast for Unsupervised Visual Representation Learning.
\newblock In \emph{CVPR}, 2020.

\bibitem{he2022mae}
K.~He, X.~Chen, S.~Xie, Y.~Li, P.~Doll{\'a}r, and R.~Girshick.
\newblock Masked Autoencoders Are Scalable Vision Learners.
\newblock In \emph{CVPR}, 2022.

\bibitem{karamcheti2023voltron}
S.~Karamcheti, S.~Nair, A.~S. Chen, T.~Kollar, C.~Finn, D.~Sadigh, and P.~Liang.
\newblock Language-Driven Representation Learning for Robotics.
\newblock In \emph{RSS}, 2023.

\bibitem{khosla2020supervised}
P.~Khosla, P.~Teterwak, C.~Wang, A.~Sarna, Y.~Tian, P.~Isola, A.~Maschinot, C.~Liu, and D.~Krishnan.
\newblock Supervised Contrastive Learning.
\newblock In \emph{NeurIPS}, 2020.

\bibitem{lerer2016learning}
A.~Lerer, S.~Gross, and R.~Fergus.
\newblock Learning Physical Intuition of Block Towers by Example.
\newblock In \emph{ICML}, 2016.

\bibitem{li2024evaluating}
X.~Li, H.~Zhang, D.~Sadigh, and D.~Xu.
\newblock Evaluating Real-World Robot Manipulation Policies in Simulation.
\newblock \emph{arXiv:2405.05941}, 2024.

\bibitem{liu2024libero}
B.~Liu, Y.~Zhu, C.~Gao, Y.~Feng, Q.~Liu, Y.~Zhu, and P.~Stone.
\newblock LIBERO: Benchmarking Knowledge Transfer for Lifelong Robot Learning.
\newblock In \emph{NeurIPS}, 2024.

\bibitem{loshchilov2019adamw}
I.~Loshchilov and F.~Hutter.
\newblock Decoupled Weight Decay Regularization.
\newblock In \emph{ICLR}, 2019.

\bibitem{ma2022vip}
Y.~J. Ma, S.~Sodhani, D.~Jayaraman, O.~Bastani, V.~Kumar, and A.~Zhang.
\newblock VIP: Towards Universal Visual Reward and Representation via Value-Implicit Pre-Training.
\newblock In \emph{ICLR}, 2023.

\bibitem{nair2022r3m}
S.~Nair, A.~Rajeswaran, V.~Kumar, C.~Finn, and A.~Gupta.
\newblock R3M: A Universal Visual Representation for Robot Manipulation.
\newblock In \emph{CoRL}, 2022.

\bibitem{oquab2024dinov2}
M.~Oquab, T.~Darcet, T.~Moutakanni, H.~Vo, M.~Szafraniec, V.~Khalidov, P.~Fernandez, D.~Haziza, F.~Massa, A.~El-Nouby, \etal.
\newblock DINOv2: Learning Robust Visual Features without Supervision.
\newblock \emph{TMLR}, 2024.

\bibitem{peng2019domainnet}
X.~Peng, Q.~Bai, X.~Xia, Z.~Huang, K.~Saenko, and B.~Wang.
\newblock Moment Matching for Multi-Source Domain Adaptation.
\newblock In \emph{ICCV}, 2019.

\bibitem{pomerleau1988alvinn}
D.~A. Pomerleau.
\newblock ALVINN: An Autonomous Land Vehicle in a Neural Network.
\newblock In \emph{NeurIPS}, 1988.

\bibitem{radford2021clip}
A.~Radford, J.~W. Kim, C.~Hallacy, A.~Ramesh, G.~Goh, S.~Agarwal, G.~Sastry, A.~Askell, P.~Mishkin, J.~Clark, \etal.
\newblock Learning Transferable Visual Models From Natural Language Supervision.
\newblock In \emph{ICML}, 2021.

\bibitem{radosavovic2023vc1}
I.~Radosavovic, T.~Xiao, S.~James, P.~Abbeel, J.~Malik, and T.~Darrell.
\newblock Real-World Robot Learning with Masked Visual Pre-training.
\newblock In \emph{CoRL}, 2023.

\bibitem{raissi2019pinn}
M.~Raissi, P.~Perdikaris, and G.~E. Karniadakis.
\newblock Physics-informed neural networks.
\newblock \emph{Journal of Computational Physics}, 2019.

\bibitem{shang2024theia}
J.~Shang, K.~Schmeckpeper, B.~B. May, M.~V. Minniti, T.~Kelestemur, D.~Watkins, and L.~Herlant.
\newblock Theia: Distilling Diverse Vision Foundation Models for Robot Learning.
\newblock In \emph{CoRL}, 2024.

\bibitem{sun2024spa}
H.~Sun, Y.~Li, J.~Zhang, Y.~Y. Schuurmans, and P.~Abbeel.
\newblock SPA: 3D Spatial-Awareness Enables Effective Embodied Representation.
\newblock \emph{arXiv:2410.08208}, 2024.

\bibitem{team2024octo}
Octo Model Team.
\newblock Octo: An Open-Source Generalist Robot Policy.
\newblock In \emph{RSS}, 2024.

\bibitem{toussaint2018differentiable}
M.~Toussaint, K.~R. Allen, K.~A. Smith, and J.~B. Tenenbaum.
\newblock Differentiable physics and stable modes for tool-use and manipulation planning.
\newblock In \emph{RSS}, 2018.

\bibitem{tschannen2025siglip2}
M.~Tschannen, A.~Gritsenko, X.~Wang, M.~S. Aga, B.~Mustafa, and N.~Houlsby.
\newblock SigLIP 2: Multilingual Vision-Language Encoders with Improved Semantic Understanding, Localization, and Dense Features.
\newblock \emph{arXiv:2502.14786}, 2025.

\bibitem{wang2024hpt}
L.~Wang, B.~T. Wen, S.~Kirmani, B.~Ichter, and D.~Sadigh.
\newblock HPT: Scaling Data and Model for Heterogeneous Pre-trained Robot Transformers.
\newblock \emph{arXiv:2409.20537}, 2024.

\bibitem{xiao2022mvp}
T.~Xiao, I.~Radosavovic, T.~Darrell, and J.~Malik.
\newblock Masked Visual Pre-training for Motor Control.
\newblock \emph{arXiv:2203.06173}, 2022.

\bibitem{yang2024physcene}
Z.~Yang, J.~Chen, M.~Li, K.~Jia, and H.~Su.
\newblock PhyScene: Physically Interactable 3D Scene Synthesis for Embodied AI.
\newblock In \emph{CVPR}, 2024.

\bibitem{ze2024generalizable}
Y.~Ze, G.~Yan, Y.-H. Wu, A.~Macaluso, Y.~Ge, J.~Ye, N.~Hansen, L.~E. Li, and X.~Wang.
\newblock Generalizable Visual Reinforcement Learning with Segment Anything Model.
\newblock \emph{arXiv:2312.17116}, 2024.

\bibitem{zhai2023siglip}
X.~Zhai, B.~Mustafa, A.~Kolesnikov, and L.~Beyer.
\newblock Sigmoid Loss for Language Image Pre-Training.
\newblock In \emph{ICCV}, 2023.

\bibitem{zhang2024physgrasp}
D.~Zhang, Z.~Wu, and S.~Song.
\newblock PhyGrasp: Generalizing Robotic Grasping with Physics-informed Large Multimodal Models.
\newblock \emph{arXiv:2402.16836}, 2024.

\bibitem{zhao2023act}
T.~Z. Zhao, V.~Kumar, S.~Levine, and C.~Finn.
\newblock Learning Fine-Grained Bimanual Manipulation with Low-Cost Hardware.
\newblock In \emph{RSS}, 2023.

\bibitem{zheng2021ressl}
M.~Zheng, S.~You, F.~Huang, C.~Wang, Z.~Su, and C.-C.~J. Kuo.
\newblock ReSSL: Relational Self-Supervised Learning with Weak Augmentation.
\newblock In \emph{NeurIPS}, 2021.

\bibitem{zhou2022ibot}
J.~Zhou, C.~Wei, H.~Wang, W.~Shen, C.~Xie, A.~Yuille, and T.~Kong.
\newblock iBOT: Image BERT Pre-training with Online Tokenizer.
\newblock In \emph{ICLR}, 2022.

\bibitem{zhu2020robosuite}
Y.~Zhu, J.~Wong, A.~Mandlekar, R.~Mart{\'i}n-Mart{\'i}n, A.~Joshi, S.~Nasiriany, and Y.~Zhu.
\newblock robosuite: A Modular Simulation Framework and Benchmark for Robot Learning.
\newblock \emph{arXiv:2009.12293}, 2020.

\end{thebibliography}
}

% ======================================================================
% Appendix
% ======================================================================
\newpage
\appendix
\section{Appendix}

\subsection{Category-Grounded Material Prior Ranges}
\label{app:priors}

Table~\ref{tab:material_priors} lists the full parameter ranges for each material type used in our category-grounded physics assignment.

\begin{table}[h]
\centering
\caption{\textbf{Material prior ranges.} Each material type defines uniform distributions over mass (kg), static friction coefficient, and coefficient of restitution.}
\label{tab:material_priors}
\small
\begin{tabular}{@{}lccc@{}}
\toprule
\textbf{Material Type} & \textbf{Mass (kg)} & \textbf{Friction} & \textbf{Restitution} \\
\midrule
Metal        & [2.0, 8.0]    & [0.10, 0.40] & [0.50, 0.90] \\
Wood         & [0.5, 3.0]    & [0.30, 0.60] & [0.30, 0.60] \\
Fabric       & [0.05, 0.50]  & [0.60, 0.95] & [0.10, 0.30] \\
Glass/Ceramic& [0.3, 3.0]    & [0.20, 0.50] & [0.50, 0.80] \\
Rubber/Plastic& [0.1, 2.0]   & [0.40, 0.80] & [0.40, 0.80] \\
Food/Organic & [0.1, 2.0]    & [0.30, 0.70] & [0.10, 0.40] \\
Paper/Light  & [0.01, 0.30]  & [0.40, 0.70] & [0.10, 0.30] \\
Animal       & [0.1, 5.0]    & [0.30, 0.60] & [0.20, 0.50] \\
Stone/Heavy  & [2.0, 15.0]   & [0.30, 0.70] & [0.10, 0.30] \\
Default      & [0.1, 5.0]    & [0.20, 0.80] & [0.10, 0.90] \\
\bottomrule
\end{tabular}
\end{table}


\subsection{LIBERO-10 Task Descriptions}
\label{app:libero_tasks}

\begin{table}[h]
\centering
\caption{\textbf{LIBERO-10 task descriptions.} All tasks are long-horizon multi-step manipulation tasks with 50 human demonstrations each.}
\label{tab:libero_tasks}
\small
\begin{tabular}{@{}cl@{}}
\toprule
\textbf{ID} & \textbf{Full Task Description} \\
\midrule
T0 & Put both the alphabet soup and the tomato sauce in the basket \\
T1 & Put both the cream cheese box and the butter in the basket \\
T2 & Turn on the stove and put the moka pot on it \\
T3 & Put the black bowl in the bottom drawer of the cabinet and close it \\
T4 & Put the white mug on the left plate and put the yellow-white mug on the right plate \\
T5 & Pick up the book and place it in the back compartment of the caddy \\
T6 & Put the white mug on the plate and put the chocolate pudding to the right of the plate \\
T7 & Put both the alphabet soup and the cream cheese box in the basket \\
T8 & Put both moka pots on the stove \\
T9 & Put the yellow and white mug in the microwave and close it \\
\bottomrule
\end{tabular}
\end{table}


\subsection{Zero-Shot $k$-NN Retrieval}
\label{app:zeroshot}

\begin{table}[h]
\centering
\caption{\textbf{Zero-shot physics retrieval.} $R^2$ for physics prediction via cosine-similarity retrieval from a 5,000-image library ($k{=}5$). Mean $R^2$ averages over all three properties.}
\label{tab:zeroshot}
\small
\begin{tabular}{@{}lcccc@{}}
\toprule
\textbf{Backbone} & \textbf{Mass $R^2$} & \textbf{Friction $R^2$} & \textbf{Rest. $R^2$} & \textbf{Mean $R^2$} \\
\midrule
SigLIP       & 0.309 & 0.199 & 0.441 & 0.316 \\
VIP          & 0.466 & 0.476 & 0.696 & 0.546 \\
R3M          & 0.487 & 0.499 & 0.681 & 0.556 \\
\methodname{} & 0.536 & 0.426 & 0.638 & 0.533 \\
CLIP         & 0.563 & \best{0.586} & 0.735 & 0.628 \\
DINOv2-B/14  & 0.582 & 0.567 & 0.741 & 0.630 \\
DINOv2-L/14  & \best{0.591} & 0.568 & \best{0.753} & \best{0.637} \\
\bottomrule
\end{tabular}
\end{table}

\subsection{Extended Training Analysis}
\label{app:extended}

Extending training from 30K to 100K steps improves clustering (NMI: $0.435 \to 0.517$, ARI: $0.335 \to 0.480$) while slightly degrading linear probing ($R^2$: $0.550 \to 0.536$ mass). We attribute this to representation compression: longer training tightens within-material clusters but smooths fine-grained physics variations exploited by linear probes. This suggests 30K steps is optimal for downstream probing, while 100K steps may be preferred for unsupervised clustering.


\subsection{LIBERO-10 Per-Seed Breakdown}
\label{app:libero_seeds}

\begin{table}[h]
\centering
\caption{\textbf{DynaCLIP LIBERO-10 per-seed results.} Success rate (\%) for each of 3 random seeds, with mean and standard deviation.}
\label{tab:libero_seeds}
\small
\begin{tabular}{@{}lcccccc@{}}
\toprule
\textbf{Task} & \textbf{Seed 0} & \textbf{Seed 1} & \textbf{Seed 2} & \textbf{Mean} & \textbf{Std} \\
\midrule
T0: Soup + sauce & 8.0 & 6.0 & 4.0 & 6.0 & 1.6 \\
T1: Cream cheese + butter & 52.0 & 62.0 & 58.0 & 57.3 & 4.1 \\
T2: Stove + moka pot & 94.0 & 92.0 & 96.0 & 94.0 & 1.6 \\
T3: Bowl in drawer & 98.0 & 96.0 & 92.0 & 95.3 & 2.5 \\
T4: Mugs on plates & 72.0 & 76.0 & 80.0 & 76.0 & 3.3 \\
T5: Book in caddy & 70.0 & 72.0 & 74.0 & 72.0 & 1.6 \\
T6: Mug + pudding & 50.0 & 58.0 & 52.0 & 53.3 & 3.4 \\
T7: Soup + cheese & 46.0 & 56.0 & 34.0 & 45.3 & 9.0 \\
T8: Both moka pots & 26.0 & 26.0 & 24.0 & 25.3 & 0.9 \\
T9: Mug in microwave & 58.0 & 66.0 & 72.0 & 65.3 & 5.7 \\
\midrule
\textbf{Average} & 57.4 & 61.0 & 58.6 & \textbf{59.0} & 2.7 \\
\bottomrule
\end{tabular}
\end{table}

\subsection{Complete LIBERO-10 Results for All Backbones}
\label{app:libero_full}

\begin{table}[h]
\centering
\caption{\textbf{Complete LIBERO-10 results.} Average success rate (\%) with standard deviation across 3 seeds for all completed backbone evaluations.}
\label{tab:libero_full}
\small
\begin{tabular}{@{}lccl@{}}
\toprule
\textbf{Backbone} & \textbf{Avg. (\%)} & \textbf{Feature Dim} & \textbf{Status} \\
\midrule
\textbf{DynaCLIP (Ours)} & \best{59.0} $\pm$ 2.7 & 1536 & Completed \\
DINOv2-B/14 & 51.5 $\pm$ 3.3 & 768 & Completed \\
CLIP-L/14 & 46.9 $\pm$ 2.9 & 768 & Completed \\
SigLIP-B/16 & \tbd{51.6} & 768 & Projected \\
Theia-B & \tbd{51.1} & 768 & Projected \\
Voltron-S & \tbd{43.9} & 384 & Projected \\
MVP & 32.9 $\pm$ 3.8 & 768 & Completed \\
MCR (MAE) & 32.9 $\pm$ 3.4 & 768 & Completed \\
VC-1-L & 31.4 $\pm$ 3.9 & 1024 & Completed \\
R3M (ResNet-50) & \tbd{15.6} & 2048 & Running \\
\bottomrule
\end{tabular}
\end{table}

\end{document}
